% ---------------------------------------------------------------------------
% Methods: the MTEX filters we compare against. Settings are those of
% src/run_mtex_baseline.m; what each filter does was checked against the
% MTEX 7.1.0 source on 2026-10-14.
% Items in [square brackets] are for you to fill in. The timings are the means
% of the mtex_log.csv files of the final Windows run (2026-10-14, 290 maps).
% ---------------------------------------------------------------------------

\subsection{Baseline filters}
\label{sec:baselines}

We compare our method with the six orientation-map filters of
MTEX~\cite{bachmann2010mtex,hielscher2019denoising}, version 7.1.0. Each
filter replaces the orientation of a pixel by one computed from the pixels
around it; they differ in how these pixels are combined.

\begin{enumerate}
  \item[(i)] \emph{Mean filter.} The pixel is replaced by the average
    orientation of the $3 \times 3$ window centred on it, that is, of itself
    and its eight neighbours. It is the simplest and fastest filter, but it
    blurs any sharp change of orientation inside the window.
  \item[(ii)] \emph{Median filter.} Of the 25 orientations in the
    $5 \times 5$ window around the pixel, the one with the smallest mean
    disorientation to all the others replaces the pixel. The result is
    always one of the measured orientations, so a few outliers in the window
    cannot pull it away.
  \item[(iii)] \emph{Kuwahara filter}~\cite{kuwahara1976processing}. The
    $5 \times 5$ window is divided into four overlapping $3 \times 3$
    quadrants, each with the pixel in one corner. The pixel takes the mean
    orientation of the quadrant in which the orientations are most alike.
    Near a sharp change this is the quadrant on the pixel's own side, so the
    change is not smeared.
  \item[(iv)] \emph{Spline filter}~\cite{garcia2010robust}. A smooth
    orientation field is fitted to the data, balancing closeness to the
    measured orientations against how strongly the field bends. The amount of
    smoothing is chosen automatically by generalised cross-validation, and a
    robust weighting reduces the influence of outliers. This is the default
    filter of MTEX.
  \item[(v)] \emph{Half-quadratic filter}~\cite{bergmann2016restoration}.
    The output map is the one that best balances closeness to the data against
    the differences between neighbouring pixels. As in total-variation
    denoising, both are measured with an $\ell_1$-type norm, which favours
    uniform regions separated by sharp steps over a gradual blur. Differences
    larger than $15^\circ$ are not penalised at all. We set the weight of the
    smoothing term to $\alpha = 0.25$, a quarter of the MTEX default.
  \item[(vi)] \emph{Infimal-convolution filter}~\cite{bergmann2018priors}.
    Like (v), but the penalty combines first-order differences (weight
    $\lambda = 0.005$), which keep steps sharp, with second-order differences
    ($\mu = 0.01$), which allow the orientation to change gradually instead
    of breaking into a staircase of uniform patches. We use the MTEX default
    parameters.
\end{enumerate}

MTEX applies a filter to one grain at a time, so that orientations are never
averaged across a grain boundary. The grains are therefore first
reconstructed from the noisy map: two neighbouring pixels belong to the same
grain when their disorientation is below a threshold angle $\omega$. Since
each pixel can be up to $\theta_{\max}$ away from its true orientation, two
pixels of the same grain can differ by up to $2\theta_{\max}$, and a smaller
threshold breaks grains into fragments. We therefore set
$\omega = \max(10^\circ, 2\theta_{\max})$, using the true $\theta_{\max}$ of
each test condition, or of each map in the mixed set. The filters are thus
given the noise level, which our network never receives.

The filters alone cannot repair misindexed pixels. A misindexed pixel has a
random orientation that almost never lies within $\omega$ of its neighbours,
so it forms a grain of its own, and MTEX leaves grains of fewer than five
pixels unfiltered. We therefore clean the map before filtering: grains
smaller than three pixels are deleted, the grains are reconstructed without
them, and each deleted pixel takes the orientation of the nearest remaining
pixel, including at the edge of the map, so that it joins that pixel's grain.
The filter is then applied to the cleaned map, grain by grain. We use the
same procedure on every map. In the seven conditions without misindexing
($p = 0$) no pixel forms a grain of its own, so the cleaning removes nothing
and the result is that of the filter alone.

The test maps were loaded into MTEX as .ang files without any change
of reference frame, which we checked reproduces every orientation exactly.
The filtered maps were written back on the original $128 \times 128$ grid and
scored in the same way as the output of our network (Eq.~\ref{eq:error}). The
filters were run on test maps 1--10 of each of the 28 test conditions and of
the mixed set, 290 maps in total. On [processor], the six filters together
took about 24~s per map, of which 16~s for the infimal-convolution filter,
against about 0.04~s per map for our network on one NVIDIA H100 GPU,
including reading and writing the files (500 maps in 21--23~s). The filter parameters were not tuned:
they are the MTEX defaults, except for the window of the median and Kuwahara
filters and $\alpha$ of the half-quadratic filter, [how these values were
chosen].
